\documentclass[12pt,a4paper]{article}
\usepackage{fontspec}
\usepackage{xeCJK}
\setCJKmainfont{Microsoft JhengHei}
\setCJKsansfont{Microsoft JhengHei}
\setCJKmonofont{Microsoft JhengHei}
\usepackage{amsmath,amssymb,amsthm}
\usepackage{geometry}
\geometry{margin=2cm}
\usepackage{enumitem}
\usepackage{fancyhdr}
\usepackage{tcolorbox}
\tcbuselibrary{theorems,skins,breakable}
\usepackage{hyperref}
\hypersetup{colorlinks=true,linkcolor=blue!70!black}

\pagestyle{fancy}
\fancyhf{}
\fancyhead[L]{\small 2022 複變函數期末考\,---\,詳解}
\fancyhead[R]{\small \thepage}
\setlength{\headheight}{14pt}
\renewcommand{\headrulewidth}{0.4pt}

\newtcolorbox{qbox}[1][]{
  colback=blue!5!white, colframe=blue!60!black,
  fonttitle=\bfseries, title=#1, breakable
}
\newtcolorbox{solbox}[1][]{
  colback=green!3!white, colframe=green!50!black,
  fonttitle=\bfseries, title=#1, breakable
}
\newtcolorbox{ansbox}{
  colback=red!5!white, colframe=red!60!black,
  fonttitle=\bfseries, title=答案, breakable
}

\newcommand{\dd}{\,\mathrm{d}}
\newcommand{\ii}{\mathrm{i}}
\newcommand{\ee}{\mathrm{e}}
\newcommand{\Res}{\operatorname{Res}}

\begin{document}

\begin{center}
{\LARGE\bfseries 2022 複變函數期末考\,---\,詳細解答}\\[6pt]
{\large 共 14 題}\\[4pt]
\end{center}

\bigskip

%==================================================================
% Q1
%==================================================================
\begin{qbox}[第 1 題]
$|z|=3$，用 ML 不等式求積分的最大值 $\displaystyle\left|\oint_C\frac{\ee^{2z}}{z+1}\dd{z}\right|=a\pi\ee^b$，求 $a-b$。
\end{qbox}

\begin{solbox}[解]
\textbf{Step 1：求 $M$}

在 $|z|=3$ 上，需要 $\left|\dfrac{\ee^{2z}}{z+1}\right|$ 的最大值。

分子：$|\ee^{2z}|=\ee^{2\text{Re}(z)}$。在 $|z|=3$ 上，$\text{Re}(z)$ 最大值為 $3$，故 $|\ee^{2z}|\le\ee^6$。

分母：$|z+1|\ge||z|-1|=|3-1|=2$，故 $\dfrac{1}{|z+1|}\le\dfrac{1}{2}$。

所以 $M=\dfrac{\ee^6}{2}$。

\textbf{Step 2：求 $L$}

$L=2\pi\cdot 3=6\pi$。

\textbf{Step 3：ML 不等式}

$$\left|\oint_C\frac{\ee^{2z}}{z+1}\dd{z}\right|\le ML=\frac{\ee^6}{2}\cdot 6\pi=3\pi\ee^6$$

故 $a=3$，$b=6$。
\end{solbox}

\begin{ansbox}
$a-b=3-6=\boxed{-3}$
\end{ansbox}

%==================================================================
% Q2
%==================================================================
\begin{qbox}[第 2 題]
計算 $\displaystyle f(z)=\oint_C\frac{4z}{(z+1)(z-1)(z-3)}\dd{z}$，

(a) $C:|z|=2$，$f(z)=a\pi\ii$，求 $a$。

(b) $C:|z|=4$，$f(z)=b\pi\ii$，求 $b$。
\end{qbox}

\begin{solbox}[解]
先做部分分式：$\dfrac{4z}{(z+1)(z-1)(z-3)}=\dfrac{A}{z+1}+\dfrac{B}{z-1}+\dfrac{C}{z-3}$

$A=\dfrac{4(-1)}{(-1-1)(-1-3)}=\dfrac{-4}{(-2)(-4)}=\dfrac{-4}{8}=-\dfrac{1}{2}$

$B=\dfrac{4(1)}{(1+1)(1-3)}=\dfrac{4}{(2)(-2)}=-1$

$C=\dfrac{4(3)}{(3+1)(3-1)}=\dfrac{12}{(4)(2)}=\dfrac{3}{2}$

\textbf{(a) $|z|=2$}：奇異點 $z=-1$ 和 $z=1$ 在圓內，$z=3$ 在圓外。

$$f(z)=\oint_C\frac{-1/2}{z+1}\dd{z}+\oint_C\frac{-1}{z-1}\dd{z}+\oint_C\frac{3/2}{z-3}\dd{z}$$
$$=-\frac{1}{2}(2\pi\ii)+(-1)(2\pi\ii)+0=-\pi\ii-2\pi\ii=-3\pi\ii$$

故 $a=-3$。

\textbf{(b) $|z|=4$}：三個奇異點 $z=-1,1,3$ 都在圓內。

$$f(z)=-\frac{1}{2}(2\pi\ii)+(-1)(2\pi\ii)+\frac{3}{2}(2\pi\ii)=-\pi\ii-2\pi\ii+3\pi\ii=0$$
\end{solbox}

\begin{ansbox}
$a=\boxed{-3}$，$b=\boxed{0}$
\end{ansbox}

%==================================================================
% Q3
%==================================================================
\begin{qbox}[第 3 題 (5-2)]
計算 $\displaystyle\int_C(x^2+\ii y)\dd{z}$，$C$ 為從 $0$ 到 $2+\ii$ 的線段。虛部值為？
\end{qbox}

\begin{solbox}[解]
$C$ 為從 $0$ 到 $2+\ii$ 的直線段，參數化：$z(t)=(2+\ii)t$，$0\le t\le 1$。

故 $x=2t$，$y=t$，$\dd{z}=(2+\ii)\dd{t}$。

$$\int_C(x^2+\ii y)\dd{z}=\int_0^1\bigl[(2t)^2+\ii t\bigr](2+\ii)\dd{t}$$
$$=\int_0^1(4t^2+\ii t)(2+\ii)\dd{t}$$
$$=\int_0^1\bigl[8t^2+4\ii t^2+2\ii t+\ii^2 t\bigr]\dd{t}$$
$$=\int_0^1\bigl[8t^2-t+\ii(4t^2+2t)\bigr]\dd{t}$$
$$=\left[\frac{8t^3}{3}-\frac{t^2}{2}\right]_0^1+\ii\left[\frac{4t^3}{3}+t^2\right]_0^1$$
$$=\left(\frac{8}{3}-\frac{1}{2}\right)+\ii\left(\frac{4}{3}+1\right)=\frac{13}{6}+\frac{7}{3}\ii$$
\end{solbox}

\begin{ansbox}
虛部 $=\boxed{\dfrac{7}{3}}$
\end{ansbox}

%==================================================================
% Q4
%==================================================================
\begin{qbox}[第 4 題 (5-5)]
$C$ 中心為原點，半徑 $2$ 的圓（$|z|=2$），$f(z)=\ee^z$，用 Cauchy's inequality 計算 $|f^{(3)}(0)|$ 的最大值。（$\ee^2\approx 7.4$）
\end{qbox}

\begin{solbox}[解]
\textbf{Cauchy 不等式}：若 $f$ 在 $|z-z_0|\le R$ 上解析，且 $|f(z)|\le M$，則
$$|f^{(n)}(z_0)|\le\frac{n!\,M}{R^n}$$

這裡 $z_0=0$，$R=2$，$n=3$。

在 $|z|=2$ 上：$|f(z)|=|\ee^z|=\ee^{\text{Re}(z)}\le\ee^2$。故 $M=\ee^2$。

$$|f^{(3)}(0)|\le\frac{3!\cdot\ee^2}{2^3}=\frac{6\ee^2}{8}=\frac{3\ee^2}{4}$$

實際上 $f^{(3)}(z)=\ee^z$，故 $f^{(3)}(0)=1$。

Cauchy 不等式給出的上界為 $\dfrac{3\ee^2}{4}\approx\dfrac{3\times 7.4}{4}=5.55$。
\end{solbox}

\begin{ansbox}
$|f^{(3)}(0)|$ 的最大值上界 $=\boxed{\dfrac{3\ee^2}{4}\approx 5.55}$
\end{ansbox}

%==================================================================
% Q5
%==================================================================
\begin{qbox}[第 5 題]
將 $f(z)=\dfrac{1}{2-z}$ 以 Maclaurin series 展開，求收斂半徑 $R$。
\end{qbox}

\begin{solbox}[解]
$$f(z)=\frac{1}{2-z}=\frac{1}{2}\cdot\frac{1}{1-\frac{z}{2}}=\frac{1}{2}\sum_{k=0}^{\infty}\left(\frac{z}{2}\right)^k=\sum_{k=0}^{\infty}\frac{z^k}{2^{k+1}}$$

收斂條件：$\left|\dfrac{z}{2}\right|<1$，即 $|z|<2$。

或用 Ratio Test：$\displaystyle\lim_{k\to\infty}\left|\frac{z^{k+1}/2^{k+2}}{z^k/2^{k+1}}\right|=\left|\frac{z}{2}\right|<1$，得 $|z|<2$。

另外，$f(z)$ 的奇異點為 $z=2$，距離原點（Maclaurin 中心）為 $2$。
\end{solbox}

\begin{ansbox}
$R=\boxed{2}$
\end{ansbox}

%==================================================================
% Q6
%==================================================================
\begin{qbox}[第 6 題]
封閉曲線 $C:|z-1|=1$，計算 $\displaystyle\oint_C\tan z\dd{z}=k\pi$，求 $k$。
\end{qbox}

\begin{solbox}[解]
\textbf{Step 1：找奇異點}

$\tan z=\dfrac{\sin z}{\cos z}$，寫成 $f(z)=\dfrac{g(z)}{h(z)}$，其中 $g(z)=\sin z$，$h(z)=\cos z$。

奇異點在 $h(z)=\cos z=0$，即 $z=\dfrac{\pi}{2}+n\pi$，$n=0,\pm 1,\pm 2,\ldots$

\textbf{Step 2：判斷哪些奇異點在 $C:|z-1|=1$ 內}

圓 $C$ 的中心為 $1$，半徑為 $1$，所以圓內區域為 $0<\text{Re}(z)<2$ 附近。

\begin{itemize}
\item $z=\dfrac{\pi}{2}\approx 1.571$：$\left|\dfrac{\pi}{2}-1\right|=\dfrac{\pi-2}{2}\approx 0.571<1$ $\Rightarrow$ \textbf{在圓內}
\item $z=-\dfrac{\pi}{2}\approx -1.571$：$\left|-\dfrac{\pi}{2}-1\right|\approx 2.571>1$ $\Rightarrow$ 在圓外
\item $z=\dfrac{3\pi}{2}\approx 4.712$：$\left|\dfrac{3\pi}{2}-1\right|\approx 3.712>1$ $\Rightarrow$ 在圓外
\end{itemize}

圓內唯一的奇異點為 $z=\dfrac{\pi}{2}$。

\textbf{Step 3：判斷極點階數}

$h(z)=\cos z$ 在 $z=\dfrac{\pi}{2}$ 的零點階數：$h\!\left(\dfrac{\pi}{2}\right)=\cos\dfrac{\pi}{2}=0$，$h'\!\left(\dfrac{\pi}{2}\right)=-\sin\dfrac{\pi}{2}=-1\ne 0$。

故 $\cos z$ 在 $z=\dfrac{\pi}{2}$ 有\textbf{一階零點}（Theorem 6.4.1），因此 $\tan z=\dfrac{\sin z}{\cos z}$ 在此有\textbf{一階極點}（Theorem 6.4.3）。

\textbf{Step 4：計算殘值}

\textbf{方法一}（Theorem 6.5.1，一階極點公式）：
$$\Res(\tan z,\tfrac{\pi}{2})=\lim_{z\to\pi/2}(z-\tfrac{\pi}{2})\cdot\frac{\sin z}{\cos z}$$

利用 L'H\^{o}pital 法則（因為分子分母在 $z=\pi/2$ 都趨近 $0$）：
$$=\lim_{z\to\pi/2}\frac{\sin z}{\frac{\cos z}{z-\pi/2}}=\lim_{z\to\pi/2}\frac{\sin z}{(\cos z)'}=\frac{\sin(\pi/2)}{-\sin(\pi/2)}=\frac{1}{-1}=-1$$

\textbf{方法二}（$g/h'$ 公式，Theorem 6.5 的替代殘值公式）：

當 $f(z)=g(z)/h(z)$，$h(z_0)=0$，$h'(z_0)\ne 0$，$g(z_0)\ne 0$ 時：
$$\Res(f(z),z_0)=\frac{g(z_0)}{h'(z_0)}=\frac{\sin(\pi/2)}{-\sin(\pi/2)}=\frac{1}{-1}=-1$$

\textbf{Step 5：由殘值定理（Theorem 6.5.3）求積分}

$$\oint_C\tan z\dd{z}=2\pi\ii\sum_k\Res(f,z_k)=2\pi\ii\cdot(-1)=-2\pi\ii$$
\end{solbox}

\begin{ansbox}
$k=\boxed{-2}$（即 $\oint_C\tan z\dd{z}=-2\pi\ii$）
\end{ansbox}

%==================================================================
% Q7
%==================================================================
\begin{qbox}[第 7 題]
$\ee^{4/z}$ 在 $z=0$ 的殘值 (residue) 為？
\end{qbox}

\begin{solbox}[解]
\textbf{Step 1：判斷奇異點類型}

$\ee^{4/z}$ 在 $z=0$ 非解析。將 $w=\dfrac{4}{z}$ 代入 $\ee^w$ 的 Maclaurin 展開（Ch 6.2），展開成 Laurent 級數。

\textbf{Step 2：Laurent 級數展開}

由 $\ee^w=\displaystyle\sum_{k=0}^{\infty}\frac{w^k}{k!}$（對所有 $w$ 收斂），令 $w=\dfrac{4}{z}$：

$$\ee^{4/z}=\sum_{k=0}^{\infty}\frac{1}{k!}\left(\frac{4}{z}\right)^k=1+\frac{4}{z}+\frac{4^2}{2!\,z^2}+\frac{4^3}{3!\,z^3}+\cdots=1+\frac{4}{z}+\frac{8}{z^2}+\frac{32}{3z^3}+\cdots$$

此級數有無窮多個負冪次項，故 $z=0$ 為\textbf{本質奇異點}（essential singularity，Ch 6.4）。

\textbf{Step 3：讀取殘值}

殘值的定義（Ch 6.5）：$\Res(f(z),z_0)=a_{-1}$，即 Laurent 級數中 $\dfrac{1}{z-z_0}$ 項的係數。

在上述展開中，$\dfrac{1}{z}$ 的係數為 $\dfrac{4^1}{1!}=4$。
\end{solbox}

\begin{ansbox}
$\Res(\ee^{4/z},0)=\boxed{4}$
\end{ansbox}

%==================================================================
% Q8
%==================================================================
\begin{qbox}[第 8 題]
將 $f(z)=\dfrac{1}{z(z-3)}$ 以 $|z|>3$ 展開成 Laurent series，求無窮級數的第三項。
\end{qbox}

\begin{solbox}[解]
部分分式：$\dfrac{1}{z(z-3)}=\dfrac{1}{z-3}\cdot\dfrac{1}{z}$

因為 $|z|>3$，即 $\left|\dfrac{3}{z}\right|<1$：

$$\frac{1}{z(z-3)}=\frac{1}{z}\cdot\frac{1}{z-3}=\frac{1}{z}\cdot\frac{1}{z(1-3/z)}=\frac{1}{z^2}\cdot\frac{1}{1-3/z}$$
$$=\frac{1}{z^2}\sum_{k=0}^{\infty}\left(\frac{3}{z}\right)^k=\sum_{k=0}^{\infty}\frac{3^k}{z^{k+2}}$$
$$=\frac{1}{z^2}+\frac{3}{z^3}+\frac{9}{z^4}+\frac{27}{z^5}+\cdots$$

第一項：$\dfrac{1}{z^2}$，第二項：$\dfrac{3}{z^3}$，第三項：$\dfrac{9}{z^4}$。
\end{solbox}

\begin{ansbox}
第三項 $=\boxed{\dfrac{9}{z^4}}$
\end{ansbox}

%==================================================================
% Q9
%==================================================================
\begin{qbox}[第 9 題]
$\displaystyle\int_0^{\infty}\frac{2\sin x}{x}\dd{x}=a+b\ii$，$a,b$ 為實數，求 $a$ 與 $b$ 的近似值。
\end{qbox}

\begin{solbox}[解]
\textbf{Step 1：建構 indented contour（Ch 6.6.2--6.6.3）}

考慮輔助函數 $f(z)=\dfrac{\ee^{\ii z}}{z}$，$z=0$ 為一階極點，位於實數軸上。

建構 indented contour $C$（避開實軸上的極點）：
\begin{itemize}
\item 實軸段 $[-R,-r]$ 和 $[r,R]$
\item $C_R$：大上半圓 $z=R\ee^{\ii\theta}$，$0\le\theta\le\pi$（逆時針）
\item $C_r$：小上半圓 $z=r\ee^{\ii\theta}$，$\theta$ 從 $\pi$ 到 $0$（順時針繞過原點）
\end{itemize}

$z=0$ 被排除在圍道外部，故 $f(z)$ 在圍道內部解析，由 Cauchy-Goursat 定理：
$$\int_{-R}^{-r}\frac{\ee^{\ii x}}{x}\dd{x}+\int_{C_r}\frac{\ee^{\ii z}}{z}\dd{z}+\int_r^R\frac{\ee^{\ii x}}{x}\dd{x}+\int_{C_R}\frac{\ee^{\ii z}}{z}\dd{z}=0 \quad\cdots(*)$$

\textbf{Step 2：大上半圓 $C_R$ 的積分（$R\to\infty$）}

$f(z)=\dfrac{\ee^{\ii z}}{z}=\dfrac{p(z)}{q(z)}\ee^{\ii z}$，其中 $p(z)=1$，$q(z)=z$。

$\deg q=1\ge\deg p+1=1$，滿足 Theorem 6.6.2 的條件，故
$$\lim_{R\to\infty}\int_{C_R}\frac{\ee^{\ii z}}{z}\dd{z}=0$$

\textbf{Step 3：小上半圓 $C_r$ 的積分（$r\to 0$）}

先求殘值：$\Res\!\left(\dfrac{\ee^{\ii z}}{z},\,0\right)=\lim_{z\to 0}z\cdot\dfrac{\ee^{\ii z}}{z}=\ee^0=1$

由 Theorem 6.6.3，對於實軸上的一階極點，小半圓積分：
$$\lim_{r\to 0}\int_{C_r}f(z)\dd{z}=\pi\ii\cdot\Res(f,0)=\pi\ii$$

但 $C_r$ 的方向是\textbf{順時針}（從 $\pi$ 到 $0$），與 Theorem 6.6.3 中定義的逆時針（$0$ 到 $\pi$）相反，故要加負號：
$$\lim_{r\to 0}\int_{C_r}\frac{\ee^{\ii z}}{z}\dd{z}=-\pi\ii$$

\textbf{Step 4：合併實軸上的積分}

對 $\displaystyle\int_{-R}^{-r}\frac{\ee^{\ii x}}{x}\dd{x}$，令 $x=-t$（$\dd{x}=-\dd{t}$）：
$$\int_{-R}^{-r}\frac{\ee^{\ii x}}{x}\dd{x}=\int_R^r\frac{\ee^{-\ii t}}{-t}(-\dd{t})=-\int_r^R\frac{\ee^{-\ii t}}{t}\dd{t}$$

與另一段合併：
$$\int_r^R\frac{\ee^{\ii x}}{x}\dd{x}-\int_r^R\frac{\ee^{-\ii x}}{x}\dd{x}=\int_r^R\frac{\ee^{\ii x}-\ee^{-\ii x}}{x}\dd{x}=\int_r^R\frac{2\ii\sin x}{x}\dd{x}$$

\textbf{Step 5：代入 $(*)$ 式取極限}

令 $R\to\infty$，$r\to 0$：
$$2\ii\int_0^{\infty}\frac{\sin x}{x}\dd{x}+(-\pi\ii)+0=0$$
$$2\ii\int_0^{\infty}\frac{\sin x}{x}\dd{x}=\pi\ii$$
$$\boxed{\int_0^{\infty}\frac{\sin x}{x}\dd{x}=\frac{\pi}{2}}$$

\textbf{Step 6：代回原題}

$$\int_0^{\infty}\frac{2\sin x}{x}\dd{x}=2\cdot\frac{\pi}{2}=\pi$$

此為實數值，故 $a=\pi\approx 3.14$，$b=0$。
\end{solbox}

\begin{ansbox}
$a\approx\boxed{3.14}$（即 $\pi$），$b=\boxed{0}$
\end{ansbox}

%==================================================================
% Q10
%==================================================================
\begin{qbox}[第 10 題]
$\displaystyle\int_0^{\pi}\frac{1}{1+2\sin^2\theta}\dd{\theta}=a+b\ii$，求 $a$ 與 $b$ 最接近的值。
\end{qbox}

\begin{solbox}[解]
\textbf{Step 1：對稱性化簡}

$\sin^2\theta$ 的週期為 $\pi$，故被積函數的週期為 $\pi$：
$$\int_0^{\pi}\frac{1}{1+2\sin^2\theta}\dd{\theta}=\frac{1}{2}\int_0^{2\pi}\frac{1}{1+2\sin^2\theta}\dd{\theta}$$

\textbf{Step 2：$z=\ee^{\ii\theta}$ 代換（Ch 6.6.1 方法）}

令 $z=\ee^{\ii\theta}$，$\dd{\theta}=\dfrac{\dd{z}}{\ii z}$，$C:|z|=1$（單位圓）。

$\sin\theta=\dfrac{z-z^{-1}}{2\ii}$，故 $\sin^2\theta=\left(\dfrac{z-z^{-1}}{2\ii}\right)^2=\dfrac{(z-z^{-1})^2}{-4}=-\dfrac{z^2-2+z^{-2}}{4}$

\textbf{Step 3：化簡被積函數}

$$1+2\sin^2\theta=1-\frac{z^2-2+z^{-2}}{2}=\frac{2-(z^2-2+z^{-2})}{2}=\frac{4-z^2-z^{-2}}{2}$$

代入積分：
$$\frac{1}{2}\oint_{|z|=1}\frac{1}{\frac{4-z^2-z^{-2}}{2}}\cdot\frac{\dd{z}}{\ii z}=\frac{1}{2}\oint_{|z|=1}\frac{2}{\ii z(4-z^2-z^{-2})}\dd{z}$$

\textbf{Step 4：消去負冪次}

將分母乘以 $z$：$z(4-z^2-z^{-2})=4z-z^3-z^{-1}=\dfrac{4z^2-z^4-1}{z}=\dfrac{-(z^4-4z^2+1)}{z}$

故
$$\frac{2}{\ii z(4-z^2-z^{-2})}=\frac{2}{\ii\cdot\frac{-(z^4-4z^2+1)}{z}}=\frac{-2z}{\ii(z^4-4z^2+1)}$$

原式 $=\dfrac{1}{2}\displaystyle\oint_{|z|=1}\frac{-2z}{\ii(z^4-4z^2+1)}\dd{z}=\frac{1}{\ii}\oint_{|z|=1}\frac{-z}{z^4-4z^2+1}\dd{z}$

即 $=\ii\displaystyle\oint_{|z|=1}\frac{z}{z^4-4z^2+1}\dd{z}$（因為 $\dfrac{-1}{\ii}=\ii$）。

\textbf{Step 5：找極點}

解 $z^4-4z^2+1=0$。令 $w=z^2$：$w^2-4w+1=0$，$w=\dfrac{4\pm\sqrt{12}}{2}=2\pm\sqrt{3}$。

\begin{itemize}
\item $z^2=2-\sqrt{3}\approx 0.268$：$z=\pm\sqrt{2-\sqrt{3}}\approx\pm 0.518$ $\Rightarrow$ $|z|\approx 0.518<1$，\textbf{在單位圓內}
\item $z^2=2+\sqrt{3}\approx 3.732$：$z=\pm\sqrt{2+\sqrt{3}}\approx\pm 1.932$ $\Rightarrow$ $|z|\approx 1.932>1$，在單位圓外
\end{itemize}

令 $\alpha=\sqrt{2-\sqrt{3}}$，則圓內有兩個一階極點 $z=\alpha$ 和 $z=-\alpha$。

分解：$z^4-4z^2+1=(z^2-(2-\sqrt{3}))(z^2-(2+\sqrt{3}))=(z-\alpha)(z+\alpha)(z^2-(2+\sqrt{3}))$

\textbf{Step 6：計算殘值}

令 $h(z)=z^4-4z^2+1$，$h'(z)=4z^3-8z$，$g(z)=z$。用 $g/h'$ 公式：

\textbf{在 $z=\alpha$}：
$$\Res\!\left(\frac{z}{z^4-4z^2+1},\,\alpha\right)=\frac{g(\alpha)}{h'(\alpha)}=\frac{\alpha}{4\alpha^3-8\alpha}=\frac{1}{4\alpha^2-8}=\frac{1}{4(2-\sqrt{3})-8}=\frac{1}{-4\sqrt{3}}$$

\textbf{在 $z=-\alpha$}：
$$\Res\!\left(\frac{z}{z^4-4z^2+1},\,-\alpha\right)=\frac{-\alpha}{4(-\alpha)^3-8(-\alpha)}=\frac{-\alpha}{-4\alpha^3+8\alpha}=\frac{-1}{-4\alpha^2+8}=\frac{-1}{4\sqrt{3}}$$

殘值總和 $=\dfrac{1}{-4\sqrt{3}}+\dfrac{-1}{4\sqrt{3}}=-\dfrac{2}{4\sqrt{3}}=-\dfrac{1}{2\sqrt{3}}$

\textbf{Step 7：殘值定理求積分}

$$\oint_{|z|=1}\frac{z}{z^4-4z^2+1}\dd{z}=2\pi\ii\cdot\left(-\frac{1}{2\sqrt{3}}\right)=-\frac{\pi\ii}{\sqrt{3}}$$

$$\int_0^{\pi}\frac{1}{1+2\sin^2\theta}\dd{\theta}=\ii\cdot\left(-\frac{\pi\ii}{\sqrt{3}}\right)=\frac{-\ii^2\pi}{\sqrt{3}}=\frac{\pi}{\sqrt{3}}$$

（利用 $\ii\cdot(-\pi\ii/\sqrt{3})=-\pi\ii^2/\sqrt{3}=\pi/\sqrt{3}$，因 $\ii^2=-1$。）
\end{solbox}

\begin{ansbox}
$a\approx\boxed{1.81}$（即 $\dfrac{\pi}{\sqrt{3}}$），$b=\boxed{0}$
\end{ansbox}

%==================================================================
% Q11
%==================================================================
\begin{qbox}[第 11 題]
$|z|=2$ 在線性分式函數 $T(z)=\dfrac{z+2}{z-2}$ 映射後為？
\end{qbox}

\begin{solbox}[解]
令 $z=2\ee^{\ii\theta}$（$|z|=2$ 上的點），代入：
$$w=T(z)=\frac{2\ee^{\ii\theta}+2}{2\ee^{\ii\theta}-2}=\frac{\ee^{\ii\theta}+1}{\ee^{\ii\theta}-1}$$

利用 $\ee^{\ii\theta}=\cos\theta+\ii\sin\theta$：

$$w=\frac{(\cos\theta+1)+\ii\sin\theta}{(\cos\theta-1)+\ii\sin\theta}$$

分子分母乘以共軛：
$$w=\frac{[(\cos\theta+1)+\ii\sin\theta][(\cos\theta-1)-\ii\sin\theta]}{|(\cos\theta-1)+\ii\sin\theta|^2}$$

分子 $=(\cos^2\theta-1)-\ii(\cos\theta+1)\sin\theta+\ii\sin\theta(\cos\theta-1)+\sin^2\theta$
$=-\sin^2\theta+\sin^2\theta+\ii\sin\theta[(\cos\theta-1)-(\cos\theta+1)]=-2\ii\sin\theta$

分母 $=(\cos\theta-1)^2+\sin^2\theta=2-2\cos\theta$。

$$w=\frac{-2\ii\sin\theta}{2(1-\cos\theta)}=\frac{-\ii\sin\theta}{1-\cos\theta}$$

$w$ 為純虛數（實部 $=0$）。因此 $|z|=2$ 映射後為\textbf{虛軸}（$\text{Re}(w)=0$）。
\end{solbox}

\begin{ansbox}
映射後為\boxed{\text{虛軸（}x=0\text{）}}
\end{ansbox}

%==================================================================
% Q12
%==================================================================
\begin{qbox}[第 12 題]
封閉曲線 $C$ 包含 $C_1$（逆時針，包含 $z=0$）和 $C_2$（順時針，包含 $z=1$）（如圖「8」字形）。

(1) $\displaystyle\oint_{C_1}\frac{2z-3}{z^2-z}\dd{z}=k_1\cdot 2\pi\ii$，求 $k_1$。

(2) $\displaystyle\oint_{C_2}\frac{2z-3}{z^2-z}\dd{z}=k_2\cdot 2\pi\ii$，求 $k_2$。
\end{qbox}

\begin{solbox}[解]
\textbf{Step 1：部分分式分解}

$$\frac{2z-3}{z^2-z}=\frac{2z-3}{z(z-1)}=\frac{A}{z}+\frac{B}{z-1}$$

cover-up method（Ch 5.3）：
\begin{itemize}
\item $A=\left.\dfrac{2z-3}{z-1}\right|_{z=0}=\dfrac{-3}{-1}=3$
\item $B=\left.\dfrac{2z-3}{z}\right|_{z=1}=\dfrac{-1}{1}=-1$
\end{itemize}

故 $\dfrac{2z-3}{z(z-1)}=\dfrac{3}{z}+\dfrac{-1}{z-1}$。

\textbf{Step 2：$C_1$（逆時針，包含 $z=0$，不包含 $z=1$）}

逐項積分（Ch 5.3，Cauchy-Goursat 定理）：
\begin{itemize}
\item $\displaystyle\oint_{C_1}\frac{3}{z}\dd{z}=3\cdot 2\pi\ii=6\pi\ii$（$z=0$ 在 $C_1$ 內，逆時針 $\Rightarrow$ $2\pi\ii$）
\item $\displaystyle\oint_{C_1}\frac{-1}{z-1}\dd{z}=0$（$z=1$ 不在 $C_1$ 內，被積函數在 $C_1$ 內解析）
\end{itemize}

$$\oint_{C_1}\frac{2z-3}{z(z-1)}\dd{z}=6\pi\ii+0=6\pi\ii=3\cdot 2\pi\ii$$

故 $k_1=3$。

\textbf{Step 3：$C_2$（順時針，包含 $z=1$，不包含 $z=0$）}

注意 $C_2$ 為\textbf{順時針}方向（負方向），$\displaystyle\oint_{\text{順}}\frac{1}{z-z_0}\dd{z}=-2\pi\ii$。

\begin{itemize}
\item $\displaystyle\oint_{C_2}\frac{3}{z}\dd{z}=0$（$z=0$ 不在 $C_2$ 內）
\item $\displaystyle\oint_{C_2}\frac{-1}{z-1}\dd{z}=(-1)\cdot(-2\pi\ii)=2\pi\ii$

（$z=1$ 在 $C_2$ 內，順時針 $\Rightarrow$ $\oint\frac{1}{z-1}\dd{z}=-2\pi\ii$，再乘以係數 $-1$）
\end{itemize}

$$\oint_{C_2}\frac{2z-3}{z(z-1)}\dd{z}=0+2\pi\ii=2\pi\ii=1\cdot 2\pi\ii$$

故 $k_2=1$。
\end{solbox}

\begin{ansbox}
$k_1=\boxed{3}$，$k_2=\boxed{1}$
\end{ansbox}

%==================================================================
% Q13
%==================================================================
\begin{qbox}[第 13 題]
封閉曲線 $C$ 為單位圓 $|z|=1$，計算

$$\oint_C\left(\frac{2}{z}+\ee^z+\sin z+3\right)\dd{z}=k\cdot 2\pi\ii$$

求 $k$。
\end{qbox}

\begin{solbox}[解]
由積分的線性性質，可逐項計算：

\textbf{(i)} $\displaystyle\oint_C\frac{2}{z}\dd{z}$：$\dfrac{2}{z}$ 在 $z=0$ 有一階極點，$z=0$ 在 $|z|=1$ 內部。

由 Cauchy 積分公式（Ch 5.5）或直接利用 $\oint_C\dfrac{1}{z}\dd{z}=2\pi\ii$：
$$\oint_C\frac{2}{z}\dd{z}=2\cdot 2\pi\ii=4\pi\ii$$

\textbf{(ii)} $\displaystyle\oint_C\ee^z\dd{z}$：$\ee^z$ 為 entire function（在整個複數平面解析），無奇異點。

由 Cauchy-Goursat 定理（Ch 5.3）：解析函數沿封閉曲線的積分為零，故 $=0$。

\textbf{(iii)} $\displaystyle\oint_C\sin z\dd{z}$：$\sin z$ 同樣為 entire function，由 Cauchy-Goursat 定理 $=0$。

\textbf{(iv)} $\displaystyle\oint_C 3\dd{z}$：常數函數為 entire function，由 Cauchy-Goursat 定理 $=0$。

合併：
$$\oint_C\left(\frac{2}{z}+\ee^z+\sin z+3\right)\dd{z}=4\pi\ii+0+0+0=4\pi\ii=2\cdot 2\pi\ii$$
\end{solbox}

\begin{ansbox}
$k=\boxed{2}$
\end{ansbox}

%==================================================================
% Q14
%==================================================================
\begin{qbox}[第 14 題]
(a) $f(z)=\dfrac{4z+6}{z-1}$，求 $\Res(f(z),1)$。

(b) $f(z)=(z+2)^2\sin\!\left(\dfrac{2}{z+2}\right)$，求 $\Res(f(z),-2)$。
\end{qbox}

\begin{solbox}[解]
\textbf{(a)} $f(z)=\dfrac{4z+6}{z-1}$，$z=1$ 為一階極點。

由 Theorem 6.5.1（一階極點殘值公式）：
$$\Res(f(z),1)=\lim_{z\to 1}(z-1)\cdot f(z)=\lim_{z\to 1}(z-1)\cdot\frac{4z+6}{z-1}=\lim_{z\to 1}(4z+6)=4(1)+6=10$$

\textbf{(b)} $f(z)=(z+2)^2\sin\!\left(\dfrac{2}{z+2}\right)$，$z=-2$ 為奇異點。

\textbf{Step 1}：令 $w=z+2$（平移至 $w=0$），則 $f=w^2\sin\!\left(\dfrac{2}{w}\right)$。

\textbf{Step 2}：利用 $\sin u$ 的 Maclaurin 展開（Ch 6.2），令 $u=\dfrac{2}{w}$：

$$\sin\!\left(\frac{2}{w}\right)=\frac{2}{w}-\frac{1}{3!}\left(\frac{2}{w}\right)^3+\frac{1}{5!}\left(\frac{2}{w}\right)^5-\cdots=\frac{2}{w}-\frac{8}{6w^3}+\frac{32}{120w^5}-\cdots$$

$$=\frac{2}{w}-\frac{4}{3w^3}+\frac{4}{15w^5}-\cdots$$

\textbf{Step 3}：乘以 $w^2$ 得 Laurent 級數：

$$w^2\sin\!\left(\frac{2}{w}\right)=w^2\cdot\frac{2}{w}-w^2\cdot\frac{4}{3w^3}+w^2\cdot\frac{4}{15w^5}-\cdots$$
$$=2w\underbrace{-\frac{4}{3}\cdot\frac{1}{w}}_{a_{-1}\text{ 項}}+\frac{4}{15w^3}-\cdots$$

\textbf{Step 4}：讀取殘值。殘值 $=a_{-1}=$ $\dfrac{1}{w}$ 項的係數 $=-\dfrac{4}{3}$。

（此為本質奇異點，因 Laurent 級數有無窮多個負冪次項。）
\end{solbox}

\begin{ansbox}
(a) $\Res(f(z),1)=\boxed{10}$

(b) $\Res(f(z),-2)=\boxed{-\dfrac{4}{3}}$
\end{ansbox}

\end{document}
